% =====================================================================
%  rectas-y-planos.tex — Unidad: Rectas y planos en el espacio
%  Requiere apuntes.sty (estilo/ en la raíz del repo; ver README).
% =====================================================================
\documentclass[11pt,a4paper]{article}
\usepackage{apuntes}

% ---------- Macros propias de esta materia (solo si hacen falta) ----------
\DeclareMathOperator{\Sol}{Sol}
\DeclareMathOperator{\dir}{dir}
\DeclareMathOperator{\dist}{dist}

% ---------- Metadatos ----------
\materia{GAL 1}
\curso{Teórico}
\unidad{Rectas y planos en el espacio}
\periodo{Clases 12 a 16 --- 2026}
\hypersetup{pdfauthor={Santiago Trujillo}}

\begin{document}
\maketitle
\vspace{-1.6em}
\begin{center}
  {\normalsize Santiago Trujillo\par}
\end{center}
\vspace{0.5em}
\tableofcontents

% =====================================================================
%  CLASES — cada clase nueva se agrega AL FINAL, sin tocar las anteriores
% =====================================================================

\clase{12}{}{El espacio de las n-uplas}

\subsection{\texorpdfstring{El espacio $\R^n$ de las $n$-uplas}{El espacio Rn de las n-uplas}}

\begin{definicion}
Dado $n\geq1$, se llama $n$-upla de reales toda sucesión finita $(x_1,\dots,x_n)$ donde $x_1,\dots,x_n\in\R$.

Se escribe $\R^n$ el conjunto de las $n$-uplas de reales.
\end{definicion}

El espacio $\R^n$ tiene la misma estructura que el espacio de las matrices columnas ($:n\times1$) o de las matrices filas ($:1\times n$).

Como para las matrices, se pueden definir:
\begin{itemize}
  \item la suma:
  \begin{align*}
  &(x_1,\dots,x_n)+(y_1,\dots,y_n)\\
  &=(x_1+y_1,\dots,x_n+y_n)
  \end{align*}
  \item la $n$-upla nula: $\vec0=\underbrace{(0,\dots,0)}_{n\text{ veces}}$

  la opuesta: $-(x_1,\dots,x_n)=(-x_1,\dots,-x_n)$
  \item el producto de un número por una $n$-upla:
  \[
  \lambda\cdot(x_1,\dots,x_n)=(\lambda x_1,\dots,\lambda x_n)
  \]
\end{itemize}

Propiedades: Para todos $u,v,w\in\R^n$.
\begin{align*}
&(u+v)+w=u+(v+w)\\
&u+v=v+u\\
&u+\vec0=\vec0+u\\
&u+(-u)=(-u)+u=\vec0\\
&0\cdot u=\vec0\\
&1\cdot u=\vec u\\
&(-1)\cdot u=-\vec u\\
&(\lambda\mu)\cdot u=\lambda(\mu\cdot u)\\
&(\lambda+\mu)\cdot u=\lambda\cdot u+\mu\cdot u\\
&\lambda(u+v)=\lambda u+\lambda v
\end{align*}
% DUDA: pág. 6 del escaneo, "u + 0 = 0 + u" queda sin "= u" al final.

Se dice que $\R^n$ es un espacio vectorial.

\subsection{\texorpdfstring{El espacio $\R^3$}{El espacio R3}}

A través de un sistema fijado de coordenadas se identifica el «espacio ambiente» con $\R^3$
\[
E=\R^3
\]

\begin{observacion}
Cada 3-upla puede representar
\begin{itemize}
  \item O bien un punto $A\in\R^3$
  \item O bien un vector $u\in\R^3$
\end{itemize}
\end{observacion}

Además las operaciones sobre $\R^3$ corresponden a las operaciones «geométricas» sobre los puntos y vectores:
\begin{itemize}
  \item $u+v$ (suma de 2 vectores)
  \item $\vec0=(0,0,0)$
  \item $-\vec u=$ % En el escaneo (pág. 7) queda sin completar.
  \item $\lambda\cdot\vec u$
  \item $A+u$
  \item $B-A=\overrightarrow{AB}$
\end{itemize}

\begin{center}
\begin{tikzpicture}[>={Stealth}]
  \draw[->] (0,0) -- (3,0.8) node[midway,above] {$v$};
  \draw[->] (0,0) -- (2,-0.9) node[midway,below] {$u$};
  \draw[->] (2,-0.9) -- (5,-0.1);
  \draw[->] (3,0.8) -- (5,-0.1);
  \draw[->] (0,0) -- (5,-0.1) node[pos=0.55,above] {$u+v$};
  \begin{scope}[xshift=7.5cm]
    \fill (0,-0.3) circle (1.5pt) node[below] {$A$};
    \draw[->] (0,-0.3) -- (2.5,-0.1) node[midway,above] {$u$};
    \node at (2.6,-0.75) {$B=A+u$};
  \end{scope}
\end{tikzpicture}
\end{center}

% =====================================================================
\clase{13}{}{Rectas en el espacio}

\subsection{\texorpdfstring{El espacio $\R^3$}{El espacio R3}}

\subsubsection*{Vectores colineales y puntos alineados}

\begin{definicion}
Dos vectores $u,v\in\R^3$ son colineales cuando existe $\lambda\in\R$ tal que $u=\lambda\cdot v$ \underline{o} $v=\lambda\cdot u$
\end{definicion}

\begin{observacion}
El «o» está para tener en cuenta el caso donde uno de los vectores es nulo:
\begin{itemize}
  \item Si $u=\vec0$ y $v\neq\vec0$: \quad $u=0\cdot v$ \quad pero \quad $v\neq\lambda\cdot u$ \quad $(\forall\lambda)$
  \item Si $u\neq\vec0$ y $v=\vec0$: \quad $v=0\cdot u$ \quad pero \quad $u\neq\lambda\cdot v$ \quad $(\forall\lambda)$
  \item Si $u,v\neq\vec0$: \quad $v=\lambda\cdot u\iff u=\lambda^{-1}\cdot v$ \quad (pues $\lambda\neq0$)
\end{itemize}
\end{observacion}

\begin{definicion}
Tres puntos $A,B,C$ son alineados cuando $\overrightarrow{AB}$ y $\overrightarrow{AC}$ son colineales.
\end{definicion}

\begin{observacion}
$A,B,C$ alineados
\end{observacion}
% En el escaneo (pág. 1) esta observación queda sin completar.

\begin{definicion}
Dados dos puntos $A,B\in\R^3$ tales que $A\neq B$, la recta $R$ que pasa por $A$ y $B$ es el conjunto de los puntos $X$ tales que $A$, $B$ y $X$ son alineados.
\[
X\in R\iff\overrightarrow{AX}\text{ y }\overrightarrow{AB}\text{ son colineales}
\]
\end{definicion}

Se puede demostrar que
\[
X\in R\iff\underbrace{X=A+\lambda(B-A)}_{\substack{\text{ecuación vectorial}\\\text{de la recta } R\\\text{(no única)}}}\quad\text{para algún }\lambda\in\R.
\]
Escribiendo $A=(a_1,a_2,a_3)$, $B=(b_1,b_2,b_3)$ y $X=(x_1,x_2,x_3)$
\[
(x_1,x_2,x_3)\in R\iff
\underbrace{\begin{cases}
x_1=a_1+\lambda(b_1-a_1)\\
x_2=a_2+\lambda(b_2-a_2)\\
x_3=a_3+\lambda(b_3-a_3)
\end{cases}}_{\substack{\text{Ecuación paramétrica de } R\\\text{(no es única)}}}
\quad\text{para algún }\lambda\in\R
\]


\begin{definicion}
Un vector director de una recta $R$ es cualquier vector $u$ de la forma $u=\overrightarrow{AB}$ $(=B-A)$ donde $A,B\in R$, $A\neq B$.
\end{definicion}

\begin{observacion}\leavevmode
\begin{enumerate}[label=(\arabic*)]
  \item Un vector director nunca es nulo
  \item todos los vectores directores de una misma recta son colineales entre sí.
  \item Si $u$ es un vector director de $R$, entonces $\lambda\cdot u$ es un vector director de $R$ para todo $\lambda\neq0$.
\end{enumerate}
\end{observacion}

\begin{observacion}
Una recta también se puede definir a partir de un punto $A$ y de un vector $u\neq0$.

(Vectorial) \quad $X=A+\lambda u$ \quad para algún $\lambda\in\R$

(Paramétrica)
$\begin{cases}
x_1=a_1+\lambda u_1\\
x_2=a_2+\lambda u_2\\
x_3=a_3+\lambda u_3
\end{cases}$
\quad para algún $\lambda\in\R$
\end{observacion}

\begin{definicion}
Dos rectas son paralelas cuando tienen los mismos vectores directores.
\end{definicion}

\begin{observacion}\leavevmode
\begin{enumerate}[label=(\arabic*)]
  \item
  \begin{align*}
  R\parallel R'&\overset{\text{def}}{\iff}\dir(R)=\dir(R')\\
  &\iff\dir(R)\cap\dir(R')\neq\emptyset
  \end{align*}
  \item Dos rectas paralelas que se encuentran en un punto son iguales.
\end{enumerate}
\end{observacion}

\begin{ejemplo}\label{ej:recta}\leavevmode
\begin{enumerate}[label=(\arabic*)]
  \item $A=(1,1,1)$ y $B=(3,5,-5)$
  \begin{align*}
  X&=(1,1,1)+\lambda(2,4,-6)\qquad(\text{para algún }\lambda)\\
  X&=(1,1,1)+\lambda(1,2,-3)
  \end{align*}
  \[
  \begin{cases}
  x_1=1+\lambda\\
  x_2=1+2\lambda\\
  x_3=1-3\lambda
  \end{cases}
  \qquad(\text{para algún }\lambda\in\R)
  \]
  \item $C=(0,-1,4)$ y $D=(1,4,-3)$ \qquad ¿$C,D\in R$?

  ¿$C\in R$?
  \[
  \begin{cases}
  0=1+\lambda\\
  -1=1+2\lambda\\
  4=1-3\lambda
  \end{cases}
  \qquad
  \begin{array}{l}
  \lambda=-1\\
  2\lambda=-2\Rightarrow\lambda=-1\\
  -3\lambda=3\Rightarrow\lambda=-1
  \end{array}
  \qquad\Sol(S)=\{-1\}
  \]
  Luego: $C\in R$.

  ¿$D\in R$?
  \[
  \begin{cases}
  1=1+\lambda\\
  4=1+2\lambda\\
  -3=1-3\lambda
  \end{cases}
  \qquad
  \begin{array}{l}
  \Rightarrow\lambda=0\\
  \Rightarrow2\lambda=3\Rightarrow\lambda=\frac32
  \end{array}
  \qquad\Sol(S)=\emptyset
  \]
  Luego: $D\notin R$
\end{enumerate}
\end{ejemplo}

\subsection{Ecuación implícita de la recta}

Sea $R$ una recta definida por un punto $A$ y un vector director $u\neq\vec0$.

Para determinar si un punto $X=(x_1,x_2,x_3)$ pertenece o no, tenemos que determinar si el sistema es compatible (en $\lambda$)
\[
\begin{cases}
x_1=a_1+\lambda u_1\\
x_2=a_2+\lambda u_2\\
x_3=a_3+\lambda u_3
\end{cases}
\]
Con matriz ampliada:
\[
\left(\begin{array}{c|c}u_1&x_1-a_1\\u_2&x_2-a_2\\u_3&x_3-a_3\end{array}\right)
\]
Se distinguen 3 casos:

\underline{Caso 1}: $u_1,u_2,u_3\neq0$
\[
\begin{array}{l}
\lambda=\dfrac{x_1-a_1}{u_1}\\[10pt]
\lambda=\dfrac{x_2-a_2}{u_2}\\[10pt]
\lambda=\dfrac{x_3-a_3}{u_3}
\end{array}
\qquad
\lambda\text{ existe si y sólo si}
\overbrace{\begin{cases}
\dfrac{x_1-a_1}{u_1}=\dfrac{x_2-a_2}{u_2}\\[10pt]
\dfrac{x_1-a_1}{u_1}=\dfrac{x_3-a_3}{u_3}
\end{cases}}^{\text{Ecuación implícita de } R}
\]
Poniendo el sistema anterior en forma lineal, se obtiene:
\[
(S)\;\begin{cases}
\dfrac{1}{u_1}x_1-\dfrac{1}{u_2}x_2=\dfrac{a_1}{u_1}-\dfrac{a_2}{u_2}\\[10pt]
\dfrac{1}{u_1}x_1-\dfrac{1}{u_3}x_3=\dfrac{a_1}{u_1}-\dfrac{a_3}{u_3}
\end{cases}
\qquad
\begin{array}{l}
\text{Sistema lineal de 2 ecuaciones}\\
\text{y 3 incógnitas } (x_1,x_2,x_3)
\end{array}
\]
Por construcción: $\Sol(S)=R$

\underline{Caso 2}: $u_1,u_2\neq0$, pero $u_3=0$
\[
\begin{cases}
x_1=a_1+\lambda u_1\\
x_2=a_2+\lambda u_2\\
x_3=a_3+\lambda u_3
\end{cases}
\qquad(\text{para algún }\lambda)
\]
\[
\begin{array}{l}
\lambda=\dfrac{x_1-a_1}{u_1}\\[10pt]
\lambda=\dfrac{x_2-a_2}{u_2}
\end{array}
\qquad
(S)\;\begin{cases}
\dfrac{x_1-a_1}{u_1}=\dfrac{x_2-a_2}{u_2}\\[10pt]
x_3=a_3
\end{cases}
\qquad\Sol(S)=R
\]

\underline{Caso 3}: $u_1\neq0$, pero $u_2=u_3=0$
\[
\begin{cases}
x_1=a_1+\lambda u_1\\
x_2=a_2\\
x_3=a_3
\end{cases}
\Longrightarrow(S):\begin{cases}x_2=a_2\\x_3=a_3\end{cases}
\]
\begin{align*}
\Sol(S)&=R\\
&=\{(x_1,a_2,a_3):x_1\in\R\}
\end{align*}

\begin{proposicion}
Toda recta tiene una ecuación implícita (o cartesiana)
\[
(S)\;\begin{cases}
a_{11}x_1+a_{12}x_2+a_{13}x_3=b_1\\
a_{21}x_1+a_{22}x_2+a_{23}x_3=b_2
\end{cases}
\]

es decir: un sistema tal que $\Sol(S)=R$
\end{proposicion}

\begin{ejemplo}[\ref{ej:recta} continuado]
$A=(1,1,1)$ y $u=(1,2,-3)$
\[
\begin{cases}
x_1=1+\lambda\Rightarrow\lambda=x_1-1\\[4pt]
x_2=1+2\lambda\Rightarrow\lambda=\dfrac{x_2-1}{2}\\[10pt]
x_3=1-3\lambda\Rightarrow\lambda=\dfrac{1-x_3}{3}
\end{cases}
\qquad
\begin{cases}
x_1-1=\dfrac{x_2-1}{2}\\[10pt]
x_1-1=\dfrac{1-x_3}{3}
\end{cases}
\]
\[
\begin{cases}
x_1-\dfrac{x_2}{2}=\dfrac12\\[10pt]
x_1+\dfrac{x_3}{3}=\dfrac43
\end{cases}
\qquad
\begin{cases}
2x_1-x_2=1\\
3x_1+x_3=4
\end{cases}
\qquad\Sol(S)=R
\]
$C=(0,-1,4)$: $C\in R$

$D=(1,4,-3)\notin R$
\end{ejemplo}

% =====================================================================
\clase{14}{}{Planos, norma y producto escalar}

\subsection{Ecuación implícita de la recta}

\begin{observacion}\leavevmode
\begin{enumerate}[label=(\arabic*)]
  \item Una ecuación implícita es un sistema lineal de 2 ecuaciones y 3 incógnitas tal que $\Sol(S)=R$. Para deducir una ecuación paramétrica de $R$, basta con resolver el sistema.
  \item Otro método para deducir una ecuación paramétrica a partir de una ecuación implícita consiste en hallar dos soluciones distintas del sistema. A partir de dichas soluciones (que son 2 puntos distintos de $R$), se construye fácilmente la ecuación paramétrica.
  \item Una ecuación implícita de una recta $R$ es la misma que una descripción de $R$ como intersección de 2 planos no paralelos.
\end{enumerate}
\end{observacion}

\subsection{Ecuación paramétrica del plano}

\begin{definicion}
Tres puntos $A,B,C\in\R^3$ \underline{no alineados} definen un plano $P$ por la ecuación vectorial:
\[
X=A+\lambda\overrightarrow{AB}+\mu\overrightarrow{AC}
\qquad
\begin{pmatrix}\text{para algunos}\\\lambda,\mu\in\R\end{pmatrix}
\]
\end{definicion}

\begin{observacion}
Como $A,B,C$ no son alineados, los vectores $\overrightarrow{AB}$ y $\overrightarrow{AC}$ no son colineales.
\end{observacion}

Ecuación paramétrica:
\[
\begin{cases}
x_1=a_1+\lambda(b_1-a_1)+\mu(c_1-a_1)\\
x_2=a_2+\lambda(b_2-a_2)+\mu(c_2-a_2)\\
x_3=a_3+\lambda(b_3-a_3)+\mu(c_3-a_3)
\end{cases}
\qquad
\begin{pmatrix}\text{para algunos}\\\lambda,\mu\in\R\end{pmatrix}
\]

\begin{definicion}
Un vector director de un plano $P$ es cualquier vector $u$ $(\neq\vec0)$ de la forma $u=\overrightarrow{AB}$, donde $A,B$ son puntos distintos de $P$.
\end{definicion}

¡Ojo! Los vectores directores de un plano no son colineales entre sí.

También se puede definir un plano a partir de un punto $A$ y dos vectores $u,v$ \underline{no colineales}:
\[
X=A+\lambda u+\mu v
\qquad
\begin{pmatrix}\text{para algunos}\\\lambda,\mu\in\R\end{pmatrix}
\]
\[
\begin{cases}
x_1=a_1+\lambda u_1+\mu v_1\\
x_2=a_2+\lambda u_2+\mu v_2\\
x_3=a_3+\lambda u_3+\mu v_3
\end{cases}
\]

\begin{definicion}
Dos planos $P$ y $P'$ son paralelos cuando tienen los mismos vectores directores:
\[
\dir(P)=\dir(P')
\]
\end{definicion}

Para verificar un plano con vectores directores $u,v\in\dir(P)$ (no colineales) es paralelo a otro plano $P'$, basta con verificar que $u,v\in\dir(P')$.

\begin{observacion}
Dos planos paralelos que se encuentran en un punto son iguales.
\end{observacion}

\subsection{Ecuación implícita de un plano}

\begin{proposicion}
Sea $P$ un plano definido por un punto $A=(a_1,a_2,a_3)$ y dos vectores directores $u=(u_1,u_2,u_3)$ y $v=(v_1,v_2,v_3)$ no colineales. Para todo $X=(x_1,x_2,x_3)$ tenemos las siguientes equivalencias.
\begin{align*}
&X\in P\\
\iff{}&\overrightarrow{AX}\text{ es combinación lineal de } u \text{ y } v\\
\iff{}&\begin{vmatrix}
x_1-a_1&u_1&v_1\\
x_2-a_2&u_2&v_2\\
x_3-a_3&u_3&v_3
\end{vmatrix}=0
\end{align*}
\end{proposicion}

Gracias a la Prop.\ se deduce que
\[
\begin{aligned}
X\in P\iff{}&(x_1-a_1)(u_2v_3-u_3v_2)\\
&-(x_2-a_2)(u_1v_3-u_3v_1)\\
&+(x_3-a_3)(u_1v_2-u_2v_1)=0\\
\iff{}&\alpha x_1+\beta x_2+\gamma x_3=\delta
\end{aligned}
\]

con
\[
\begin{aligned}
\alpha&=(u_2v_3-u_3v_2)\\
\beta&=(u_3v_1-u_1v_3)\\
\gamma&=(u_1v_2-u_2v_1)
\end{aligned}
\qquad\qquad
\delta=\det(A,u,v)
\]

Recíprocamente, toda ecuación de la forma
\[
\alpha x_1+\beta x_2+\gamma x_3=\delta
\]
donde $(\alpha,\beta,\gamma)\neq(0,0,0)$ define un plano.

\begin{observacion}
Para deducir la ecuación paramétrica a partir de la implícita basta con resolver el sistema (con 1 ecuación y 3 incógnitas)
\end{observacion}

\subsection{Norma y producto escalar}

\begin{definicion}
Se define la norma $\norm{u}$ de un vector $u\in\R^3$ por
\[
\norm{u}=\dist(A,B)\quad\text{donde } A,B \text{ son tales que } u=\overrightarrow{AB}
\]
\end{definicion}

\begin{proposicion}\leavevmode
\begin{enumerate}[label=(\arabic*)]
  \item $\norm{u}\geq0$ y $\norm{u}=0$ si y sólo si $u=\vec0$
  \item $\norm{\lambda\cdot u}=\abs{\lambda}\cdot\norm{u}$
  \item $\norm{u+v}\leq\norm{u}+\norm{v}$ \quad desigualdad triangular.
\end{enumerate}
\end{proposicion}

\begin{definicion}
Un versor es un vector de norma 1
\end{definicion}

\begin{definicion}
Dados dos vectores $u,v\in\R^3$, se define el \underline{producto escalar} (producto interno) por
\[
u\cdot v:=\norm{u}\cdot\norm{v}\cdot\cos\varphi
\]
donde $\varphi$ es el ángulo formado por los vectores $u$ y $v$.

Cuando $u$ o $v$ es nulo, \quad $u\cdot v=0$
\end{definicion}

\begin{notacion}
$u\cdot v=\langle u,v\rangle=\langle u\,|\,v\rangle$
\end{notacion}

\subsubsection*{Propiedades del producto escalar}

\begin{enumerate}[label=(\arabic*)]
  \item $u\cdot v=v\cdot u$
  \item $(\lambda\cdot u)\cdot v=u\cdot(\lambda\cdot v)=\lambda\cdot(u\cdot v)$
  \item $(u+u')\cdot v=u\cdot v+u'\cdot v$

  $u\cdot(v+v')=u\cdot v+u\cdot v'$
\end{enumerate}

% =====================================================================
\clase{15}{}{Producto escalar y vectorial}

\subsection{Norma y producto escalar}

\begin{definicion}
La norma de un vector $u$ en $\R^3$ es el número $\norm{u}$ definido por
\[
\norm{u}=\dist(A,B),\quad\text{donde } A,B\in\R^3 \text{ tal que } u=\overrightarrow{AB}=B-A
\]
\end{definicion}

\begin{observacion}
La def no depende de $A$ y $B$
\end{observacion}

\begin{proposicion}\leavevmode
\begin{enumerate}[label=(\arabic*)]
  \item $\norm{u}\geq0$ y $\norm{u}=0$ si y sólo si $u=\vec0$
  \item $\norm{\lambda\cdot u}=\abs{\lambda}\cdot\norm{u}$
  \item $\norm{u+v}\leq\norm{u}+\norm{v}$
\end{enumerate}
\end{proposicion}

\begin{definicion}
Un versor es un vector $u$ tal que $\norm{u}=1$
\end{definicion}

Si $u\in\R^3$ es tal que $u\neq\vec0$, entonces el vector $\dfrac{1}{\norm{u}}\cdot u$ es un versor

\begin{definicion}
El producto escalar de dos vectores $u,v\in\R^3$ es el número $u\cdot v$ $\bigl(\langle u,v\rangle,\ \langle u\,|\,v\rangle\bigr)$ definido por:
\[
u\cdot v=\norm{u}\cdot\norm{v}\cdot\cos\varphi,\quad\text{donde }\varphi\text{ es el ángulo entre } u \text{ y } v.
\]
Cuando $u$ o $v$ es nulo: \quad $u\cdot v=0$
\end{definicion}

\begin{center}
\begin{tikzpicture}[>={Stealth},scale=1.1]
  % u y v
  \draw[->,thick] (0,0) -- (2,0) node[midway,below] {$u$};
  \draw[->,thick] (0,0) -- (2.6,1.8) node[pos=0.8,above left] {$v$};
  \draw (0.6,0) arc (0:34.7:0.6); \node at (0.85,0.25) {$\varphi$};
  % proyeccion de v sobre la recta de u
  \draw[dashed] (-0.5,0) -- (0,0); \draw[dashed] (2,0) -- (3.3,0);
  \draw[dashed] (2.6,1.8) -- (2.6,-0.9);
  \draw (2.45,0) -- (2.45,0.15) -- (2.6,0.15);
  % v'
  \draw[->] (0,-0.7) -- (2.6,-0.7) node[midway,below] {$v'$};
  \draw[dashed] (0,0) -- (0,-0.9);
  % proyeccion de u sobre la recta de v
  \draw[dashed] (2,0) -- (1.352,0.936);
  \draw[dashed] (2.6,1.8) -- (3.1,2.15); \draw[dashed] (-0.3,-0.21) -- (0,0);
  % u'
  \draw[->] (-0.34,0.49) -- (1.01,1.43) node[midway,above left] {$u'$};
  \draw[dashed] (0,0) -- (-0.34,0.49);
  \draw[dashed] (1.352,0.936) -- (1.01,1.43);
\end{tikzpicture}
\end{center}
\begin{align*}
u\cdot v&=\norm{u}\underbrace{\bigl(\norm{v}\cdot\cos\varphi\bigr)}_{\norm{v'}}=\norm{u}\cdot\norm{v'}\\
&=\norm{u'}\cdot\norm{v}
\end{align*}

\begin{proposicion}\leavevmode
\begin{enumerate}[label=(\arabic*)]
  \item $u\cdot v=v\cdot u$
  \item $(\lambda u)\cdot v=u\cdot(\lambda v)=\lambda(u\cdot v)$
  \item $(u_1+u_2)\cdot v=u_1\cdot v+u_2\cdot v$

  $u\cdot(v_1+v_2)=u\cdot v_1+u\cdot v_2$
\end{enumerate}
\end{proposicion}

\begin{definicion}
Dos vectores $u$ y $v$ son ortogonales cuando $u\cdot v=0$
\end{definicion}

\subsection{Expresión en una base ortonormal}

El espacio «ambiente» se ha identificado con $\R^3$ a través de un sistema de coordenadas $(O,\vec\imath,\vec\jmath,\vec k)$.

Ahora se supone que el sistema de coordenadas $(O,\vec\imath,\vec\jmath,\vec k)$ es ortonormal, es decir
\begin{enumerate}[label=(\arabic*)]
  \item $\norm{\vec\imath}=\norm{\vec\jmath}=\norm{\vec k}=1$ \quad (son versores)
  \item $\vec\imath\cdot\vec\jmath=\vec\imath\cdot\vec k=\vec\jmath\cdot\vec k=0$ \quad (son ortogonales dos a dos)
\end{enumerate}

\begin{proposicion}
En una base ortonormal, para todo $u=(u_1,u_2,u_3)$ y $v=(v_1,v_2,v_3)$, tenemos que:
\begin{enumerate}[label=(\arabic*)]
  \item $u\cdot v=u_1v_1+u_2v_2+u_3v_3$
  \item $\norm{u}=\sqrt{u\cdot u}=\sqrt{u_1^2+u_2^2+u_3^2}$
\end{enumerate}
\end{proposicion}

\begin{demo}
\begin{align*}
\vec u&=u_1\cdot\vec\imath+u_2\cdot\vec\jmath+u_3\cdot\vec k\\
\vec v&=v_1\cdot\vec\imath+v_2\cdot\vec\jmath+v_3\cdot\vec k
\end{align*}
\begin{align*}
u\cdot v={}&(u_1\cdot\vec\imath+u_2\cdot\vec\jmath+u_3\cdot\vec k)\cdot(v_1\cdot\vec\imath+v_2\cdot\vec\jmath+v_3\cdot\vec k)\\
={}&\boxed{u_1v_1(\vec\imath\cdot\vec\imath)}+u_1v_2(\vec\imath\cdot\vec\jmath)+u_1v_3(\vec\imath\cdot\vec k)+{}\\
&u_2v_1(\vec\jmath\cdot\vec\imath)+\boxed{u_2v_2(\vec\jmath\cdot\vec\jmath)}+u_2v_3(\vec\jmath\cdot\vec k)+{}\\
&u_3v_1(\vec k\cdot\vec\imath)+u_3v_2(\vec k\cdot\vec\jmath)+\boxed{u_3v_3(\vec k\cdot\vec k)}\\
={}&u_1v_1+u_2v_2+u_3v_3
\end{align*}
\end{demo}

\begin{demo}[Demostración de (2)]
\begin{align*}
u\cdot u&=\norm{u}\cdot\norm{u}\cdot\cos0\\
&=\norm{u}^2
\end{align*}
\end{demo}

Moraleja: Cuando se consideran productos escalares, normas, distancias y ángulos, siempre se trabaja en una base ortonormal. En este marco:
\[
u\cdot v=u_1v_1+u_2v_2+u_3v_3=\begin{pmatrix}u_1&u_2&u_3\end{pmatrix}\cdot\begin{pmatrix}v_1\\v_2\\v_3\end{pmatrix}
\]
A partir del producto escalar, se deducen:
\begin{itemize}
  \item[] (normas) \quad $\norm{u}=\sqrt{u\cdot u}=\sqrt{u_1^2+u_2^2+u_3^2}$
  \item[] (distancias) \quad $AB=\norm{\overrightarrow{AB}}=\norm{B-A}=\sqrt{(b_1-a_1)^2+(b_2-a_2)^2+(b_3-a_3)^2}$
  \item[] (ángulos) \quad $\cos\varphi=\dfrac{u\cdot v}{\norm{u}\cdot\norm{v}}$ \quad donde $u,v\neq\vec0$, $\varphi$ es el ángulo entre $u$ y $v$
\end{itemize}

\subsection{Producto vectorial}

\begin{definicion}
Sea $(\vec\imath,\vec\jmath,\vec k)$ una base (ortonormal)
\begin{center}
\begin{tikzpicture}[>={Stealth}]
  \fill (0,0) circle (1.3pt) node[left] {$\vec 0$};
  \draw[->] (0,0) -- (0,1.5) node[right] {$\vec k$};
  \draw[->] (0,0) -- (1.2,-0.5) node[below] {$\vec\jmath$};
  \draw[->] (0,0) -- (-0.9,-0.7) node[below] {$\vec\imath$};
  \node[align=center] at (0,-1.6) {sentido antihorario\\positivo};
  \begin{scope}[xshift=6cm]
    \draw[->] (0,0) -- (0,1.5) node[right] {$\vec k$};
    \draw[->] (0,0) -- (1.2,-0.5) node[below] {$\vec\imath$};
    \draw[->] (0,0) -- (-0.9,-0.7) node[below] {$\vec\jmath$};
    \node[align=center] at (0,-1.6) {sentido horario\\negativo};
  \end{scope}
\end{tikzpicture}
\end{center}
\end{definicion}

\begin{definicion}
El producto vectorial de dos vectores $u,v\in\R^3$ es el vector $u\wedge v$ definido por:
\begin{enumerate}[label=(\arabic*)]
  \item $\norm{u\wedge v}=\norm{u}\cdot\norm{v}\cdot\sen\varphi$ \quad donde $\varphi$ es el ángulo entre $u$ y $v$
  \item $(u\wedge v)\cdot u=(u\wedge v)\cdot v=0$

  ($u\wedge v$ es ortogonal a $u$ y a $v$)
  \item la terna $(u,v,u\wedge v)$ es positiva
\end{enumerate}
\end{definicion}

Propiedades:
\begin{enumerate}[label=(\arabic*)]
  \item $\norm{u\wedge v}$ es el doble del área del triángulo formado por $u$ y $v$
  \item $u\wedge v=\vec0$ si y sólo si $u$ y $v$ son colineales
  \item El producto vectorial es \underline{anticonmutativo}
  \[
  u\wedge v=-(v\wedge u)
  \]
  \item $(\lambda u)\wedge v=u\wedge(\lambda v)=\lambda(u\wedge v)$
  \item $(u_1+u_2)\wedge v=(u_1\wedge v)+(u_2\wedge v)$

  $u\wedge(v_1+v_2)=(u\wedge v_1)+(u\wedge v_2)$
\end{enumerate}

¡Ojo! «$\wedge$» no es asociativa:
\[
u\wedge(v\wedge w)\neq(u\wedge v)\wedge w
\]
Pero:
\begin{enumerate}[label=(\arabic*),start=6]
  \item $(u\wedge v)\wedge w+(v\wedge w)\wedge u+(w\wedge u)\wedge v=\vec0$
\end{enumerate}

\subsection{Expresión en una base ortonormal positiva}

\begin{proposicion}
Si $(\vec\imath,\vec\jmath,\vec k)$ es una base ortonormal positiva, entonces para todos $u,v\in\R^3$
\[
u\wedge v=\begin{pmatrix}
u_2v_3-u_3v_2\\
u_3v_1-u_1v_3\\
u_1v_2-u_2v_1
\end{pmatrix}
\]
\end{proposicion}

% =====================================================================
\clase{16}{}{Normal, ángulos y distancias}

\subsection{Aplicaciones a la geometría}

Los productos $u\cdot v$ y $u\wedge w$ permiten calcular distancias y ángulos:
\begin{align*}
d(A,B)&=\norm{\overrightarrow{AB}}=\sqrt{\overrightarrow{AB}\cdot\overrightarrow{AB}}\\
&=\sqrt{(b_1-a_1)^2+(b_2-a_2)^2+(b_3-a_3)^2}
\end{align*}
\[
\cos\varphi=\frac{u\cdot v}{\norm{u}\cdot\norm{v}}
\]
\begin{center}
$\displaystyle\sen\varphi=\frac{\norm{u\wedge v}}{\norm{u}\cdot\norm{v}}$
\end{center}

\subsection{Normal a un plano}

En lo que sigue, se trabaja en una base ortonormal positiva.

\begin{definicion}
Un vector $n\neq0$ es normal a un plano $P$ cuando es ortogonal a todo vector director de $P$.
\end{definicion}

En particular, si $P$ está definido por un punto $A$ y 2 vectores directores $u,v$ no colineales:
\[
n\text{ es normal a } P\iff n\cdot u=n\cdot v=0
\]

\begin{observacion}\leavevmode
\begin{enumerate}[label=(\arabic*)]
  \item todos los vectores normales de un plano son colineales entre sí.
  \item Si $u,v$ son vectores directores no colineales de $P$, entonces el vector $u\wedge v$ es normal a $P$.
\end{enumerate}
\end{observacion}

También se puede definir un plano $P$ por un punto y un vector normal $n$:

Para todo $X\in\R^3$, tenemos que
\begin{align*}
X\in P&\iff\overrightarrow{AX}\perp n\\
&\iff\overrightarrow{AX}\cdot n=0\\
&\iff(x_1-a_1)n_1+(x_2-a_2)n_2+(x_3-a_3)n_3=0\\
&\iff\alpha x_1+\beta x_2+\gamma x_3+\delta=0
\end{align*}
donde $\alpha=n_1$, $\beta=n_2$, $\gamma=n_3$
\[
\delta=-A\cdot n
\]

\underline{Recíprocamente}: Si $P$ tiene ecuación
\[
\alpha x_1+\beta x_2+\gamma x_3+\delta=0
\]
entonces $n:=(\alpha,\beta,\gamma)$ es un vector normal de $P$

\subsubsection*{Paramétrica \texorpdfstring{$\to$}{->} Implícita}

\[
X=A+\lambda u+\mu v
\qquad
\begin{array}{l}
(u,v\text{ no colineales})\\
(\lambda,\mu\in\R)
\end{array}
\]
Sea $n:=u\wedge v=(\alpha,\beta,\gamma)$
\[
\alpha x_1+\beta x_2+\gamma x_3+\delta=0
\]
donde $\delta=-A\cdot n$

\subsubsection*{Implícita \texorpdfstring{$\to$}{->} Paramétrica}

\[
P:\ \alpha x_1+\beta x_2+\gamma x_3+\delta=0
\]
Sean: \quad $n=(\alpha,\beta,\gamma)$ vector normal

$A$ un punto de $P$ (hallar una solución)

Sea $u:=n\wedge(\text{cualquier vector no colineal a } n)$

\qquad o bien $\vec\imath$, $\vec\jmath$ o $\vec k$
\[
v:=n\wedge u
\]
\[
P:\ X=A+\lambda u+\mu v
\]

\subsection{Ángulos entre rectas y planos}

\[
\begin{aligned}
R&:\ X=A+\lambda u\\
R'&:\ X=A'+\lambda u'
\end{aligned}
\]
\begin{center}
$\displaystyle\cos\varphi=\frac{u\cdot u'}{\norm{u}\cdot\norm{u'}}$
\end{center}

\[
\begin{aligned}
P&:\ \alpha x_1+\beta x_2+\gamma x_3+\delta=0\\
P'&:\ \alpha'x_1+\beta'x_2+\gamma'x_3+\delta'=0
\end{aligned}
\]
\[
\cos\varphi=\frac{n\cdot n'}{\norm{n}\cdot\norm{n'}}
\qquad
\begin{aligned}
\text{donde }n&=(\alpha,\beta,\gamma)\\
n'&=(\alpha',\beta',\gamma')
\end{aligned}
\]

\[
\begin{aligned}
R&:\ X=A+\lambda u\\
P'&:\ \alpha'x_1+\beta'x_2+\gamma'x_3+\delta'=0
\end{aligned}
\]
\[
\cos\Bigl(\frac{\pi}{2}-\varphi\Bigr)=\frac{u\cdot n'}{\norm{u}\cdot\norm{n'}}
\]

\subsection{\texorpdfstring{Distancia de un punto $Q$ a un plano $P$}{Distancia de un punto Q a un plano P}}

\[
\dist(\underset{\mathclap{\substack{\uparrow\\\text{Punto}}}}{Q},\underset{\mathclap{\substack{\uparrow\\\text{Plano}}}}{P})=\min\{\dist(Q,X):X\in P\}
\]

\begin{itemize}
  \item $P$ plano definido por un punto $A$ y un vector normal $n$
  \item $Q$ punto cualquiera
\end{itemize}
\[
\dist(Q,P)=\frac{\abs{\overrightarrow{AQ}\cdot\vec n}}{\norm{\vec n}}
=\frac{\abs{\alpha q_1+\beta q_2+\gamma q_3+\delta}}{\sqrt{\alpha^2+\beta^2+\gamma^2}}
\]

% =====================================================================
\end{document}
